\mxnewpage\section*{Correction EV4-B}\label{EV4-B-CORR}
\textbf{QCM :} 1 A, 2 B, 3 C, 4 C, 5 D, 6 B, 7 C, 8 C, 9 D, 10 C, 11 A, 12 D. Chaque réponse correcte : \num{0,5} point ; total 6.
\paragraph{Vérifications des automatismes.} Priorités ; somme de fractions ; soustraction des exposants ; coefficient multiplicateur \num{1,2} ; division par \num{0,8} ; produit \num{1,25} fois \num{0,8} ; résolution affine ; accroissement vertical divisé par l’horizontal ; moyenne ; angle divisé par 360 ; complément à 1 ; substitution entre parenthèses.
\subsection*{Exercice 1 -- Un tarif moyen}
\begin{enumerate}
\item Valeurs 0,4,6 ; probabilités \num{0,3} ; \num{0,3} ; \num{0,4}. \textit{Barème : 1 point. Résultat et justification élémentaire attendus ; demi-point possible pour démarche correcte inachevée.}
\item $E(X)=\num{3,6}=\dfrac{18}{5}$ ; $E(X^2)=\num{19,2}=\dfrac{96}{5}$ ; variance $\dfrac{156}{25}$. \textit{Barème : 2 points. Méthode ou justification : 1 ; calcul et conclusion : 1.}
\item Moyenne $\dfrac{13}{5}$ ; écart type $\dfrac{2\sqrt {39}}{5}$, inchangé. Le prix -1 serait un remboursement : vérifier si cette politique est réellement prévue. \textit{Barème : 2 points. Méthode ou justification : 1 ; calcul et conclusion : 1.}
\end{enumerate}
\subsection*{Exercice 2 -- Maximum exponentiel}
\begin{enumerate}
\item $f\prime=t(2-t)e^{-t}$. \textit{Barème : 1 point. Résultat et justification élémentaire attendus ; demi-point possible pour démarche correcte inachevée.}
\item Croît de 0 à 2 puis décroît ; maximum $4e^{-2}$. \textit{Barème : 2 points. Méthode ou justification : 1 ; calcul et conclusion : 1.}
\item $f(4)=16e^{-4}$ ; rapport au maximum $4e^{-2}<1$ puisque $e^2>4$. \textit{Barème : 2 points. Méthode ou justification : 1 ; calcul et conclusion : 1.}
\end{enumerate}
\subsection*{Exercice 3 -- Déterminer une position}
\begin{enumerate}
\item $\pt{-x;-2}\cdot\pt{4-x;-2}=x^2-4x+4=(x-2)^2$. \textit{Barème : 1 point. Résultat et justification élémentaire attendus ; demi-point possible pour démarche correcte inachevée.}
\item Produit nul seulement en x=2. \textit{Barème : 1 point. Résultat et justification élémentaire attendus ; demi-point possible pour démarche correcte inachevée.}
\item Cercle de centre (2;0), rayon 2 ; la droite y=2 lui est tangente en (2;2). \textit{Barème : 2 points. Méthode ou justification : 1 ; calcul et conclusion : 1.}
\end{enumerate}
\paragraph{Étalonnage.} Durée cible : 20–25 min pour les automatismes, 80–85 min pour les exercices, 10–15 min de relecture. Ces durées sont des estimations éditoriales à mesurer en classe. Accepter toute méthode correcte ; ne pas sanctionner deux fois une erreur numérique propagée avec cohérence. Le barème proposé ne vaut pas décision d’harmonisation de l’équipe.
